Przygotować się na pytanie o testy statystyczne.

Przygotować się do obrony: dlaczego eksperyment A i eksperyment B testowano osobno, bez wariantu łączonego.

Sprawdzić jak cytować repozytorium we fragmencie pracy magisterskiej.

Dla pewności przejrzeć wymagania:
KRYTERIUM

PRACA MAGISTERSKA

Tytuł i temat pracy

Tytuł pracy dyplomowej powinien być jednoznaczny, zwięzły, jasno określony i adekwatny do treści. Temat pracy dotyczy wyraźnie wyodrębnionego problemu lub zagadnienia, którego rozwiązania lub opracowania podejmuje się autor. Problematyka pracy jest zgodna z kierunkiem studiów.

Charakter pracy

Może mieć charakter:

(1)   Analityczno-badawczy poprzez uzyskanie nowych wyników badań, ich jakościową i ilościową analizę, ich interpretację, wykrywanie nowych mechanizmów i zależności, nowych aspektów zjawisk i procesów.

(2)   Analityczno-projektowy poprzez zaproponowanie nowych rozwiązań lub usprawnień w nawiązaniu do aktualnego stanu wiedzy naukowo-technicznej.

(3)   Przeglądowy poprzez usystematyzowane przedstawienie oraz krytyczne omówienie zagadnień teoretycznych i technicznych dotyczących danego wycinka rzeczywistości w świetle aktualnej literatury naukowej.*

* Praca o charakterze przeglądowym jest dopuszczalna na kierunkach studiów, których efekty uczenia się dopuszczają taką formę.

Cel pracy

W pracy formułuje się cele o charakterze badawczym wymagające doboru i zastosowania metod badawczych, wykorzystując wiedzę teoretyczną oraz naukową.

Wskazane jest przedstawienie, co nowego jest zaproponowane w pracy oraz podanie ograniczeń i słabych/mocnych stron opracowanego rozwiązania (jeżeli dotyczy).

Praca ma odpowiedzieć na pytanie, czy poziom wiedzy i umiejętności autora predysponuje go do rozwiązywania problemów badawczych.

Praca jest kompletnym pod względem merytorycznym opracowaniem, potwierdzającym umiejętność samodzielnego rozwiązania problemu badawczego.

Struktura i zawartość pracy

Ma postać pisemnego monograficznego opracowania, podzielonego na rozdziały i podrozdziały. Wszystkie części pracy są powiązane z realizacją celu pracy.

W strukturze pracy wyraźnie wyodrębniono:

·         Wstęp zawierający cele pracy, motywacje jej podjęcia oraz krótki opis jej zawartości.

·         Część literaturową nawiązującą do tematu i problematyki pracy, w której – w świetle istniejącej literatury naukowej i specjalistycznej – scharakteryzowane są podstawowe terminy, aktualny dorobek, teoretyczne ujęcia i stan dotychczasowych wyników badań dotyczących badanego zjawiska.

·         Część badawczą, w której sformułowane są cele badań, problemy i pytania badawcze, a także opisane zastosowane metody, techniki i narzędzia badawcze. Ponadto w tej części pracy powinien się znaleźć opis przebiegu badań lub studium przypadku wraz z interpretacją uzyskanych wyników.

·         Zakończenie, które podsumowuje najważniejsze wnioski, podaje możliwości dalszego rozwinięcia wykonanych prac i wskazuje obszar potencjalnego zastosowania pracy.

·         Bibliografię zawierającą aktualne, wyczerpujące i wiarygodne źródła, w tym naukowe – co najmniej kilkanaście pozycji, z przewagą artykułów naukowych, ale tylko te prace, na które autor powołuje się w tekście. W przypadku pracy o charakterze przeglądowym bibliografia powinna zawierać co najmniej kilkadziesiąt artykułów naukowych z czasopism lub materiałów konferencyjnych.

Rezultaty pracy

Rezultaty pracy mają charakter poznawczy, mogą mieć charakter użytkowy. Należy dokonać analizy uzyskanych wyników. Rezultaty powinny charakteryzować się oryginalnością, a nawet w pewnym stopniu nowatorstwem.

Redakcja pracy

Praca jest przygotowana z zachowaniem zalecanych wymogów edytorskich przyjętych na Wydziale. Praca jest napisana poprawnie stylistycznie i gramatycznie, bez używania kolokwializmów i języka żargonowego. Język pracy jest rzeczowy, zwięzły i klarowny. W pracy używa się poprawnej terminologii, właściwej dla informatyki. Dodatkowe elementy pracy, jak tabele, wykresy, rysunki itp. są czytelne i dobrze ilustrują omawiane treści.

) czy wszystko jest spełnione
"na potrzeby pracy badano hipoteze - odnieść się do tytułu, detekcja wykorzystujące głbokie sieci i mały reżim danych
